\section{Motion Estimation}
Motion estimation finds a model that relates image content across frames. Its output predicts where pixels or blocks moved and is used to reduce temporal redundancy in video coding.
\subsection{Variational Methods}
Variational motion estimation defines motion as the field minimizing an energy functional. The objective balances agreement with observed frames against regularity assumptions such as spatial smoothness.
\subsubsection{Velocity Vector Field - Optical Flow}
Optical flow is a dense apparent-velocity field assigning a horizontal and vertical velocity to each image position. Velocity is measured in pixels per unit time, while displacement over one frame interval is measured in pixels.
$$V : (x,y) \in \mathcal{J} \subset \mathbb{R}^2 \to (u,v)$$
First part of the process is the formulation, then there is the discretization part.
\subsubsection{Motion Vector Field}
A motion vector field is the discrete displacement assigned to pixels or blocks between reference and current frames. In video coding it is side information, so its accuracy must be balanced against the bits required to transmit it.
$$\forall \text{ point } \Longrightarrow \exists \text{ vector that represents pixel velocity of motion}$$
Formulation:
$$\arg\min_{(u,v)} \left\{D\left[f_0, f_1, (u,v)\right] + R(u,v)\right\} \qquad \text{ where}\begin{cases}
    D \equiv \text{ Data-attachment Term}\\
    R \equiv \text{ Regularization Term}
\end{cases}$$
\subsubsection{Data Attachment and Regularization}
\begin{itemize}
    \item \textbf{Data Attachment}: measures coherence of the estimated field compared to $f_0, f_1$\\
    Each vector should link pixels of the same color = same pixel.\\
    Problem: we have much more pixels than colors (assuming $2^8$ colors).
    \item \textbf{Regularization} $R$: We include A-priori knowledge on the MVF.\\
    It penalizes MVF which are different than our A-priori knowledge, it doesn't depend on data.
\end{itemize}
\subsubsection{Optical Flow Problem}
Optical Flow $\propto$ color of the pixels: the goal is to find the best possible predictor. Problem Statement:
\begin{itemize}
    \item Motion = estimated on the Y component of YCbCr
    \item we have a Generic Point $\vec{p} = \left[x, y\right]^T$
    \item we have a trajectory $x(t)$
\end{itemize}
The optical flow is the estimate velocity of trajectory of point $\vec{p}$ at time $t$.\\
The associated displacement is the MVF.
\subsubsection{Optical Flow - Displacement Field}
Consider at a given time $t_0 \to x(t_0) = \vec{p} - {D}$\\
Then for a time step $T \to x(t_0 + T) = \vec{p}$\\
We can define the displacement as this difference:
$$D(\vec{p}, t_0, T) = x(t_0 + T) - x(t_0) = \left[\begin{matrix}
    c(\vec{p},t_0, T)\\
    d(\vec{p},t_0, T)
\end{matrix}\right] = \left[\begin{matrix}
    c(x,y)\\
    d(x,y)
\end{matrix}\right] \qquad \text{ where }\begin{cases}
    c = \text{ Horizontal Displacement}\\
    d = \text{ Vertical Displacement}
\end{cases}$$
\subsubsection{Optical Flow - Constant Illumination Hypothesis (CIH)}
CIH is one of the \textbf{Fundamental Hypothesis}.
\begin{center}
    \textit{`Luminance doesn't change along the motion trajectory $x(t)$'}
\end{center}
Mathematically we can use the brightness $f$ to define it:
$$f(x,y, t+T) = f\left(x-c(x,y),y-d(x,y),t\right)$$
However this hypothesis is rarely satisfied, due to discretization (sampling, aliasing), noise, degradation of signal.\\
$c,d$ are unknowns, we can estimate them by linearizing $f$, we assume $f$ is linear w.r.t. its own data.
\subsubsection{Optical Flow - Equation}
By using Taylor's expansion:
$$f(\vec{p},t+T) = f(\vec{p},t) - \left[c(\vec{p})\frac{\partial f(\vec{p},t)}{\partial x} + d(\vec{p})\frac{\partial f(\vec{p},t)}{\partial y}\right] + o\left[\left\|D(\vec{p})\right\|\right]$$
By dividing all by T and rearranging terms:
$$\frac{f(\vec{p},t+T) -f(\vec{p},t)}{T}= -\frac{1}{T}\left[c(\vec{p})\frac{\partial f(\vec{p},t)}{\partial x} + d(\vec{p})\frac{\partial f(\vec{p},t)}{\partial y}\right] +\frac{1}{T}o\left[\left\|D(\vec{p})\right\|\right]= -\frac{\vec{D}\nabla f}{T} + \frac{1}{T}o\left[\left\|D(\vec{p})\right\|\right]$$
We can see the left-side term is the incremental ratio (derivative), then from CIH:
$$\frac{\partial f}{\partial t} = -\vec{V}\nabla f \qquad \Longrightarrow \qquad u\frac{\partial f}{\partial x} + v\frac{\partial f}{\partial y} + \frac{\partial f}{\partial t} = \boxed{uf_x + vf_y + f_t = 0}$$
\subsubsection{Optical Flow - Solution: Horn \& Schunck method}
How to find $u$ and $v$? We use the energy of the function.\\
Assuming the total variation of the velocity over a region $\mathcal{R}$ is:
$$\iint_{\mathcal{R}}(uf_x + vf_y +f_t)^2 dx dy = {\min} \text{ s.t. }\iint \left(\left\|\nabla u\right\|^2 + \left\|\nabla v\right\|^2\right) dx dy \leq \tau$$
It is a combined minimization problem, it can be solved using lagrangian multipliers:
$$J = \iint_{\mathcal{R}}(uf_x + vf_y +f_t)^2 dx dy + \lambda\left[\iint \left(\left\|\nabla u\right\|^2 + \left\|\nabla v\right\|^2\right) dx dy - \tau\right]$$
$$\text{Solutions }\qquad \longrightarrow \begin{cases}
    \lambda \nabla_2 u = (uf_x + vf_y +f_t)f_x\\
    \lambda \nabla_2 v = (uf_x + vf_y +f_t)f_y
\end{cases} \text{ and with }\begin{cases}
    \nabla u = \bar{u} - u\\
    \nabla v = \bar{v} - v
\end{cases} \longrightarrow \begin{cases}
    u = \bar{u} - f_x \frac{\bar{u}f_x + \bar{v}f_y + f_t}{\lambda +\left\|\nabla f\right\|^2}\\
    v = \bar{v} - f_y \frac{\bar{u}f_x + \bar{v}f_y + f_t}{\lambda +\left\|\nabla f\right\|^2}
\end{cases}$$
Horn and Schunck method provides a smooth motion estimation (smooth and coherent result), very used in practice for object tracking, video compression and motion analysis.
\subsection{Block Matching}
Block matching estimates motion by searching a reference frame for the block most similar to each current block. It produces one displacement vector per block rather than per pixel, reducing complexity and signaling rate at the cost of a coarser motion model.
We split images into blocks $B_{p,q}$ with $(p,q) \in N \times M$\\
We can estimate our global motion using simple motions of small blocks, simple to program.\\
The idea is to compare a single block luminance $f_k(B_{p,q})$ with a subblocks in a block window $f_h(B_{p-i, q-j})$
\subsubsection{Formulation}
Imagine we have two blocks in position $B_{p,q}$ and $B_{p-i,q-j}$, then image indexes $h,k$ and window $W$ (full or subset of the image). We can find the luminance values and compare them using the function $d$:
$$d\left(f_k(B_{p,q}),f_h(B_{p-i,q-j})\right) = J(i,j) \qquad \Longrightarrow \qquad (\hat{i},\hat{j}) = \arg\min_{(i,j) \in W}J(i,j)$$
Different techniques differ in:
\begin{itemize}
    \item minimization criterion of $J$
    \item set of candidate vectors $W$
    \item Block size and shape
\end{itemize}
Basically there is a function assigning a score depending on different parameters: we try to minimize that score.
\subsubsection{Evaluation of the MVF}
How can we tell if the result we get is good enough?
\begin{itemize}
    \item \textbf{Motion-Compensation (MC) prediction associated to MVF} (pretty similar to the other equations) $$\text{Similar to CIH - Brightness }\tilde{f}_k(n,m) = f_h(n + u_{h\to k}(n,m), m + v_{h\to k}(n,m))$$
    \item \textbf{Prediction Error} $$e(n,n) = f_k(n,m) - \tilde{f}_k(n,m)$$
    \item \textbf{MC-ed MSE} $$ \mathcal{E} = \frac{1}{NM}\sr{}{n,m}e^2(n,m)$$
    \item \textbf{PSNR} $=10\log_{10}\frac{255^2}{\mathcal{E}}$
\end{itemize}
Main Tradeoff on complexity is to change block size $P \times Q$.\\
Larger blocks reduce coding costs, complexity, and increase MSE
\subsubsection{Block Matching criteria}
The matching criterion is a distortion score used to rank candidate reference blocks. Lower SAD or SSD means a closer pixel match; a regularized criterion also charges the estimated number of bits needed to code the motion vector.
\begin{itemize}
    \item \textbf{SSD} (Sum of Squared Differences), norm-based criteria $$J_{SSD}(i,j) = \sr{}{(n,m)\in B_{p,q}}\left[f(n,m,k) - f(n-i, m-j, h)\right]^2$$
    \begin{itemize}
        \item PROS: Good for maximizing PSNR
        \item CONS: Difficult to compute, increase in entropy of MVF, doesn't account global illumination changes.
    \end{itemize}
    \item \textbf{SAD} (Sum of Absolute Differences) $$J_{SAD}(i,j) = \sr{}{(n,m)\in B_{p,q}}\left|f(n,m,k) - f(n-i, m-j, h)\right|$$
    \begin{itemize}
        \item PROS: globally costs less bits, quality still really good
        \item CONS: problems with parts with same colors (unknown behavior due to noise)
    \end{itemize}
    \item \textbf{REG} Regularized norm-based criteria $\to$ solves most of the problems $$J_{REG}(i,j) = \left\|\vec{f}_k(B_{p,q}) - \vec{f}_h(B_{p-i,q-j})\right\|^p_p + \lambda R(i,j)$$
\end{itemize}
